\section{Open socio-technical problems in mechanistic interpretability}
\label{sec:open-problems-sociotechnical}

Effective practical application of mechanistic interpretability brings both technical challenges and complex social ramifications. It could enable us to act on AI policy and governance, presenting a valuable opportunity to implement regulatory standards and social ideals through technical means (\Cref{subsec:technical-progress}). Such consequential impacts inevitably come with important social and philosophical considerations, which likewise require rigorous inquiry if we are to fully realize the potential benefits of AI (\Cref{subsec:social-philosophical}).


\subsection{Translating technical progress in mechanistic interpretability into levers for AI policy and governance}
\label{subsec:technical-progress}

Current frontier AI governance efforts rarely specify concrete ways in which a mechanistic understanding of AI models might be used to help implementation. For example, OpenAI's Preparedness Framework commits to mitigating biological risks posed by its AI systems \citep{openai2023preparedness}, but the framework lacks details on specific measures that might be taken. Progress in interpretability could potentially enable the removal of any knowledge from the model which could aid users in creating biological weapons \citep{li2024wmdp}. However, it remains uncertain how technical progress will translate into better AI governance, largely due to the numerous open technical problems in mechanistic interpretability. However, there are several promising routes toward better levers for AI policy and governance.

These avenues include assisting companies and governments to identify risks through evaluations and enhancing forecasts about new AI developments; more thorough oversight of AI systems in deployment; simplifying how AI systems operate within existing liability law through clearer explanations of AI decisions; enabling governments to establish risk mitigation regulations and companies to commit to concrete mitigation commitments; and protecting copyright law.

An understanding of model internals can help AI labs assess the risks from frontier models \citep{Chang_2024_SurveyLLM, shevlane2023modelevaluationextremerisks,Casper_2024} and thus better fulfill their obligations under the EU AI Act to “perform model evaluation ... with a view to identifying and mitigating systemic risk” \citep{eu2024ai}. More specifically, a mechanistic understanding of models could help evaluators elicit dangerous capabilities via improved finetuning \citep{ukAISI2024early}, guide their adversarial red-team attempts \citep{tong2024massproducingfailuresmultimodalsystems}, and ensure that AI systems are not intentionally underperforming evaluations \citep{vanderweij2024aisandbagginglanguagemodels}. Advancements in mechanistic interpretability could also assist companies and governments in anticipating when or if AI models will obtain specific dangerous capabilities. Improved forecasting capabilities could enhance threat modeling by reducing the “reasonable disagreement amongst experts over which risks to prioritize” \citep{anthropic_rsp}. For instance, it may help build consensus on whether language models are just stochastic parrots \citep{bender2021dangers} or if they have coherent world models \citep{li2024emergentworldrepresentationsexploring}. We also might be able to use interpretability to build evidence supporting or refuting different models of catastrophic threat, such as determining the validity of mesa-optimization \citep{hubinger2021riskslearnedoptimizationadvanced} or inner alignment concerns \citep{carlsmith2023schemingaisaisfake}. It would also give additional time for companies to prepare adequate risk mitigation measures, \citep{openai2023preparedness}, and for governments to establish appropriate guidance or regulation \citep{ukgov2022ai}.

The EU AI Act mandates that developers of General Purpose AI models with systemic risk have obligations to report incidents involving their system to the AI office. Interpretability tools have the potential to continuously monitor AI inference and detect incidents that require reporting. Compared to incidents in other domains (for instance, nuclear security), AI systems allow us to log all inputs and system states, even those that may lead to catastrophic harm. Access to a small number of such data points may greatly improve our ability to mitigate similar future failures \citep{greenblatt2024catching}. For example, interpretability could be used to investigate the critical “features” of the input that led to the AI incident. This could improve our ability to further red-team the system (\Cref{subsec:better-predictions}) and generate more similar data points that could result in similar incidents. This, in turn, may help us reduce the likelihood of future incidents \citep{chan2024visibilityaiagents}. The incident could be utilized more directly in the construction of test-time monitors to detect similar future incidents (see \citet{roger2023coup}). More speculatively, interpretability could be used to verify companies' compliance with domestic regulation \citep{obrien2023deploymentcorrectionsincidentresponse}, or to authenticate states' compliance with future international treaties regarding the use of AI \citep{aarne2024secure}.

Leveraging a mechanistic understanding of model internals could also make decision rationales for AI model outputs more easily obtainable. This could aid in enforcing citizens' rights under the EU General Data Protection Regulation “to obtain an explanation of the decision reached” by a system “based solely on automated processing” \citep{GDPR_2016, gilpin2019explainingexplanationssociety}. Model editing tools could also resolve problems regarding the copyright status of existing generative models \citep{grynbaum2023times}. According to the US Copyright Act \citep{USCopyright2022}, for any copyrighted work, an artifact “from which the work can be perceived, reproduced, or otherwise communicated…with the aid of a machine or device” is considered a copy of the work \citep{lee2024talkinboutaigeneration}. Interpretability tools could help detect and remove memorized works that can be reproduced verbatim by generative models.


\subsection{Social and philosophical problems in mechanistic interpretability}
\label{subsec:social-philosophical}

The ability to interpret advanced AI systems holds immense potential to advance the science of AI and increase our ability to control it \citep{critch2020airesearchconsiderationshuman,tegmark2023provablysafesystemspath,dalrymple2024guaranteedsafeaiframework}. In this paper, we provide an overview of various interpretability tools that offer novel insights. Nonetheless, interpretability research has thus far produced few tools that are used to make state-of-the-art systems safer in the real world \citep{rauker2023transparentaisurveyinterpreting}. Modern AI systems are still generally trained, evaluated, monitored, and debugged using techniques that do not rely on understanding their internal workings.

The absence of paradigmatic clarity is a major socio-technical factor for this. Questions such as which goals the field of interpretability should pursue, how success should be graded, and how we should define interpretability warrant more thoughtful answers than exist at present. In interpretable AI research, the motivations and methods employed are often described as “diverse and occasionally discordant” \citep{Lipton_2016_InterpretabilityMotivations}. At times, the objective of AI interpretability research is articulated as advancing a fundamental “understanding” or “uncovering the true essence” \citep{christensen2015peering} of what is happening inside black-box models. The intrinsic validity of this paradigm deserves philosophical inquiry. 

However, from an engineer's perspective, pursuing “understanding” without a practical downstream application misses an engineer's objective. Proponents of this view may contend that quantifiable benchmarks linked to concrete practical goals are accurate measures of the success of interpretability. This motivation is concrete and useful, but some have criticized interpretability research as artificially limiting the solution space to engineering problems. When interpretability tools are studied with motivations such as fairness or safety, \citet{krishnan2020against} argues that, “Since addressing these problems need not involve something that looks like an `interpretation' (etc.) of an algorithm, the focus on interpretability artificially constrains the solution space by characterizing one possible solution as the problem itself.” Thus, some argue that interpretability research has failed to produce competitive techniques \citep{rauker2023transparentaisurveyinterpreting,Caspar_2023_RedTeamingDNNFeatureSynthesis}, omitted non-interpretability baselines \citep{rudin2019stop, krishnan2020against}, and graded interpretability tools on their own curve \citep{doshivelez2017rigorousscienceinterpretablemachine, MILLER_2019_Explanation, rauker2023transparentaisurveyinterpreting}. In all safety-relevant applications of mechanistic interpretability, it is important to assess the usefulness of interpretability against alternative methodologies. Failing to do so or misrepresenting these comparisons can lead to follow-up work that rests on false assumptions \citep{lipton2018troublingtrendsmachinelearning, leech2024questionablepracticesmachinelearning}.This is particularly problematic when it impacts the efforts of critical safety work.


Despite these concerns, it is not always imperative for mechanistic interpretability to strictly outperform uninterpretable baselines: It may still be helpful to develop methods that offer interpretability-based advantages, along with the benefits of more competitive uninterpretable methods (e.g. \citet{panickssery2024steeringllama2contrastive}). While interpretability-based methods do not currently outperform black-box baselines, if they perform in a similar ballpark, further progress in mechanistic interpretability could soon lead to better methods, especially in problems which we believe might disproportionately benefit from improved mechanistic understanding (\Cref{sec:open-problems-applications}).

Another potential reason for lack of clarity is that models are often studied in a vacuum: \citet{Smart_2024} emphasize that the usefulness or correctness of model interpretations can depend on the broader context of the model's development or deployment. For example, it may be hard to identify representations of fairness within models, since understanding this requires an understanding of broader contexts. The same data may lead to different conclusions under different definitions of fairness.

Finally, the interpretability community must exercise caution in their communication to minimize potential abuses of the results of its work. Unfortunately, selective transparency can be used to actively mislead \citep{Ananny_2018_LimitationsTransparency}. Furthermore, interpretability is at risk of being used for purposes that might serve corporate interests at the potential expense of safety. The field of AI interpretability is highly influenced by research teams in the private sector. On one hand, industry resources and research contributions have significantly advanced interpretability research. On the other hand, compared to academia, corporations publish selectively, often have financial conflicts of interest, and may provide limited transparency. Meanwhile, the recent definite — but ultimately modest — progress in mechanistic interpretability has been used to lobby against specific AI regulation by falsely claiming that, “Although advocates for AI safety guidelines often allude to the `black box' nature of AI models, where the logic behind their conclusions is not transparent, recent advancements in the AI sector have resolved this issue, thereby ensuring the integrity of open-source code models.” \citep{a16z_2023_lobbying}.

